\chapter{来源路径与阅读边界}
\label{app:sources}

三国的政治年序、县廷事务、边疆关系与书写文化保存在不同文类中。阅读一个事件时，应先问材料最接近什么，再问它没有保存谁的声音；来源不是替叙事背书的标签，而是决定一句话能说到哪里的路径。

\section{正史与编年：把并行的年序放在一起}
\label{src:chronicles}

\sourcebook{三国志}魏、蜀、吴诸书与\sourcebook{资治通鉴}是共同年序的主要入口。它们使\eventref{TK-0222-YILING}{夷陵}、\eventref{TK-0249-GAOPING-TOMBS}{高平陵}和\eventref{TK-0263-CHENGDU-SURRENDER}{成都出降}可以放进相互牵动的时间中；魏纪年的框架、篇幅的不对称和后出的编年重组，也要求读者不要把任一政权的记录当作全景。

\sourcebook{后汉书·献帝纪}与\sourcebook{晋书}帝纪、列传尤其用于核对禅代、受降与伐吴的顺序。它们能照亮\eventref{TK-0265-WEI-ABDICATION}{魏晋禅代}和\eventref{TK-0280-SUN-HAO-SURRENDER}{孙皓出降}的仪式性步骤，不能单独复原诏书起草、群臣选择或地方接收。

\section{文书与物质材料：看见不同尺度的国家}
\label{src:documents}

走马楼吴简、正始三体石经残石、城址与仓储遗存不替代国史，却使县廷登记、官学书写和物质条件变得可见。它们改变我们阅读\eventref{TK-0240-ZHENGSHI-STONE-CLASSICS}{正始三体石经}与\eventref{TK-0268-TAISHI-CODE}{《泰始律》}的尺度：一块石经或一部律文可以证明文本、规则和公共空间，不能证明各地已经同样实行。

\section{区域材料：让道路与边缘反驳中心视角}
\label{src:regional}

\sourcebook{华阳国志}、裴松之注所引材料、区域史、山城与道路考古，为南中、巴蜀、辽东、高句丽与海上问题提供必要的比较。它们帮助说明\eventref{TK-0225-SHU-NANZHONG}{南征}、\eventref{TK-0238-SIMA-YI-LIAODONG}{平辽东}和\eventref{TK-0238-HIMIKO-EMBASSY}{魏受倭使}所处的空间，不把后出的地名比定、中央战报或族名误作精确疆界与地方社会的全貌。

\section{怎样沿本卷复核}
\label{src:routes}

可从\hyperref[app:index:events]{事件索引}回到首次叙事，再由\hyperref[app:index:themes]{政权与主题索引}比较不同国家的问题；\hyperref[app:index:figures]{图件索引}说明所涉走廊与图示限度，\hyperref[app:debates]{分歧附录}集中处理会改变解释的争议。这样阅读，来源路径服务于事件，而不会把读者带回抽象的材料清单。
